\begin{landscape}
\begin{table}[htbp]
\centering
\begin{threeparttable}
\caption{Literature Positioning: Comparison of Empirical Designs}
\label{tab:literature_comparison}
\small
\begin{tabular}{llllll}
\toprule
Paper & Data & Horizon & Main Expectation Object & Identification & Mechanism Direct? \\
\midrule
\citet{Bordalo2020b} & SPF professional forecasts & Quarterly & Revision $\to$ error predictability & Cross-forecaster variation & No (inferred) \\
\citet{Coibion2015} & SPF, Michigan, Livingston & Quarterly & Revision $\to$ error dynamics & Time-series moments & No (inferred) \\
\citet{Cavallo2017} & Online survey experiment & 3--6 months & Belief updating after info treatment & Randomized information & Yes (experimental) \\
\citet{DacuntoHoangPirolliWeber2021} & German household panel & Monthly & Grocery exposure $\to$ expectations & Price scanner data & Yes (direct exposure) \\
\citet{Malmendier2016} & Michigan Survey & Quarterly & Lifetime experience $\to$ expectations & Cohort variation & No (inferred) \\
Andre et al.\ (2022) & Survey experiment & N/A & Subjective macro models & Randomized scenarios & Yes (direct elicitation) \\
\midrule
\textbf{This paper} & \textbf{CFPS + PBoC} & \textbf{Biennial / Quarterly} & \textbf{Revision $\to$ error (household)} & \textbf{Province $\times$ wave FE} & \textbf{No (inferred)} \\
\bottomrule
\end{tabular}
\begin{tablenotes}[flushleft]
\footnotesize
\item \textit{Notes:} This table compares the empirical design of the present paper with related work on household expectation formation. ``Mechanism Direct?'' indicates whether the paper directly measures the hypothesized belief-formation channel or infers it from reduced-form moments.
\end{tablenotes}
\end{threeparttable}
\end{table}
\end{landscape}
